% ৩.১৭ এনকোডার
\section{Encoder}

\subsection{সংজ্ঞা ও কাজ}
\begin{itemize}
\item Encoder কি: 2\textsuperscript{n} input থেকে n-bit output
\item Priority encoder ধারণা
\end{itemize}

\subsection{ব্যবহার}
\begin{itemize}
\item Keyboard encoding, data compression
\end{itemize}

\begin{learningbox}
এই টপিক শেষে শিক্ষার্থী পারবে:
\begin{itemize}
\item encoder-এর input-output relation লিখতে।
\item $2^n$ input থেকে $n$ output ধারণা বুঝতে।
\item priority encoder-এর প্রয়োজনীয়তা ব্যাখ্যা করতে।
\end{itemize}
\end{learningbox}

\begin{definitionbox}{Encoder}
যে combinational circuit একাধিক input line-এর সক্রিয় অবস্থাকে কম সংখ্যক binary output line-এ code করে, তাকে encoder বলা হয়। সাধারণ encoder-এ $2^n$ input এবং $n$ output থাকে।
\end{definitionbox}

\begin{center}
\fbox{\parbox{0.88\textwidth}{
\centering
\textbf{Block diagram placeholder: 4-to-2 Encoder}\\[4pt]
Inputs: $D_0,D_1,D_2,D_3$ $\rightarrow$ Encoder block $\rightarrow$ Outputs: $Y_1,Y_0$
}}
\end{center}

\begin{center}
\begin{tabular}{|c|c|c|}
\hline
\textbf{Active input} & $Y_1$ & $Y_0$ \\
\hline
$D_0$ & 0 & 0 \\
\hline
$D_1$ & 0 & 1 \\
\hline
$D_2$ & 1 & 0 \\
\hline
$D_3$ & 1 & 1 \\
\hline
\end{tabular}
\end{center}

\begin{mistakebox}{সাধারণ ভুল}
Encoder-এ সাধারণত এক সময়ে একটি input active ধরা হয়। একাধিক input active হলে priority encoder দরকার হয়।
\end{mistakebox}
